\documentclass[12pt,a4paper]{article}
\usepackage[utf8]{inputenc}
\usepackage[T1]{fontenc}
\usepackage{polski}
\usepackage[polish]{babel}
\usepackage{geometry}
\geometry{margin=2cm}
\usepackage{graphicx}
\usepackage{tabularx}
\usepackage{booktabs} 
\usepackage{hyperref}
\usepackage{bookmark}
\usepackage{enumitem}
\usepackage{titlesec}
\usepackage{microtype}
\usepackage{varwidth}
\usepackage{multicol}


\title{
    \begin{center}
    \includegraphics[width=\textwidth]{fig/logo.png}
    \vspace{0.25cm}
    \end{center}
    \textbf{\huge Bazy danych} \\
    \vspace{0.5cm}
    \textit{Projekt: GG Pwr }
    }
\author{Piotr Frąc, Oskar Pilczuk}
\date{Październik 2026}


\begin{document}
\maketitle
\clearpage


\begin{multicols*}{2}
\subsection*{Wprowadzenie}
GG PWr - \textbf{Giełda Grup PWr} - 

\subsection*{Zakres projektu}
\noindent\textbf{Musi mieć}
\begin{itemize}
    \item System rejestracji i logowania
    \item Tworzenie, edycja i usuwanie postów o grupach
    \item Możliwość filtrowania 
    \item Pobrania kalenadrza zajęć!
    \item Walidacja 
    \item Usuwanie postów po danym terminie
    \item Chat user
    \item Podział na role
\end{itemize}

\noindent\textbf{Powinien mieć}
\begin{itemize}
    \item Resotowanie hasła i mfa
    \item Template maila do porwadzących
\end{itemize}

\noindent\textbf{Mogłoby być}
\begin{itemize}
    \item Auth Google i Usos
\end{itemize}


\noindent\textbf{Nie będzie miał}
\begin{itemize}
    \item Hashowania wiadomości
\end{itemize}

\end{multicols*}


\end{document}